% !TeX root = ../papers/nc-p3.tex
\section{全ての中心鉛筆を許した曲線分岐の排除}
\label{sec:unrestricted-corridor-exclusion}
\begin{theorem}\label{thm:unrestricted-corridor-exclusion}
全射影の次元が1以下の標準正則窓には順・逆いずれの分岐中心もない。
\end{theorem}
\begin{proof}[次数・中心型の有限な範囲]
分岐があるとして、各位相の順・逆中心数を $F_i,B_i$ とする。
定理\ref{thm:finite-pure-relation-locus}と命題\ref{prop:finite-pure-pencil-opposite}より
各々高々2で、各中心に主曲線次数 $e=0,1,2,3$ と例外方向高々 $3-e$ がある。$e=0$ は主像が定点の場合である。
同位相同方向に二中心があれば両方 $e>0$。両向きで少なくとも二位相が占有される
（定理\ref{thm:single-pure-branch-exclusion}の第四次数の下界）。
順中心と逆中心を実際の列により全単射で対応させる。中心から $m$ 個の動く曲線を経て終点に至るので
終点位相は始点位相 $i$ に $m+1$ を足したもの。位相ごとの全次数和は7以下である。
ある初回移動の位相では、それ以外の全移動は主曲線一つ以下または例外直線に着く。
各成分を二度使えず、これらの総数は高々3である（$e=0$ なら例外だけ）。
よって同位相の移動は高々4回、$2\le m\le20$ である。

各列の $d_1=d_m=1$、$d_2=3-f_1$ と逆側の同式、および
命題\ref{prop:branch-corridor-degree}の三階差を用いる。$m>2$ なら $d_2,d_{m-1}\ge2$。
初回以外の移動を各該当中心の主像または例外直線に割り当てる。別の中心の初回なら主像である。
動く曲線からの移動と別に、次位相の初回直線外の中心は各々一つの例外直線を占める。
その数は少なくとも $F_{i+1}-f_1$ であり、これらの直線は一般前点が固定なので移動先と重複しない。
$d_2=2$ なら命題\ref{prop:conic-corridor-excess}の $A(t)$ により初回直線上の中心の $e$ は $d_3$ である。
非初回の移動でその前の写像が埋込みなら、始点曲線の次数は主像へ行く場合は2以下、
例外直線へ行く場合は1以下である（定理\ref{thm:sparse-pure-edge-exclusion}の射影の証明）。
ここで埋込みと数えるのは次数1,2、初回と終回から二番目、または主像の有理正規三次曲線である。

始点の中心 $p$ の初回直線 $L$ と、位相 $i+2$ の逆中心 $z$ の前点直線 $M$ を比較する。
交わらなければ $M$ からの $L$ 中心の射影は全単射である。$M$ では後点が常に $z$ なので
$e>0$ の主曲線と両立せず、$e=0$ でも $z$ はただ一つの固定主像に限られる。
交点があればその後点 $z$ は $p$ の第二曲線にあり、逆欠損に少なくとも1寄与する。
よって $b_2\ge B_{i+2}-\mathbf1_{e=0}$ を得る。逆向きにも同じである。
各順欠損の上界は $3d_j$ で、順中心のない位相では0である。
埋込み曲線上で、その位相の全順中心が二直線列の始点で各列のいずれかの端の $e$ が2以上なら、
上界を $F_{i+j}$ に下げられる。実際過剰接触から命題\ref{prop:excess-contact-two-line-pure}で得る
二関係は両端の縮約三空間に属し、両者に長さ2以上の純関係を与えて $3-e\le1$ に反する。
逆欠損についても同様である。
\end{proof}

\begin{proof}[配置の列挙と矛盾]
同位相の初回直線は互いに交わらない。従って $f_1>0$ の列の本数は次位相の順中心数以下である。
$f_2>0$ の列も同じ二位相先の順中心を共有できない。共有すればその一意な前点が
二本の初回直線に属すからである。$f_2=3d_2-3-d_3$ であり、二直線列の場合は $d_3=0$ とする。
同位相に二直線列が二本あれば各始点で $e>0$、かつその終回直線には逆欠損2がある。
上の直線交叉の議論から、この位相の全逆中心が両方の終回直線に属してしまう。
終回直線も互いに交わらないので不可能である。これらを逆向きにも使う。

初回二番目の三次曲線と、同位相に来る終回から二番目の別の三次曲線も比較する。
両始終端の $e$ が3、両列の長さが4以上なら、それぞれを定める二つの移動元は
他方の初回直線と異なるので、互いに相手の主三次曲線である。
同じ列なら二つの成分の一致は列内の再訪となる。別列でも一般前点の一意性により成分は異なる。
定理\ref{thm:branch-cubic-contact}の単元による小行列式の包含はそのまま使える。
前者の逆欠損、後者の順欠損は互いに素な台で両曲線の交叉に含まれるので、その和は5以下である。
両側の次数列を $d,r$（後者は終点から逆順）と書くと、その和の下界は
$(d_4-3d_3+8)+(r_4-3r_3+8)$ である。

以上の整数条件の全列挙は付属コードで行う。まず総数の等しい二組の位相数を全て取り、
各始終点・長さについて許容次数を小さい添字から追加し、確定した三階差の符号を照合する。
次に相異なる終点への完全マッチングを作り、位相次数和と各中心の移動先の本数を加算する。
各 $e$ を全て割り当て、三階差が欠損の上下界の差に入るかを検査し、上記の交叉を検査する。
全列が $(1,3,1)$ の場合は相対二位相各一中心か各二中心に限られ、前者は
定理\ref{thm:single-pure-branch-exclusion}末尾の二補完式、後者は全中心 $e=3$ なので
命題\ref{prop:branch-local-pure-pencil}で消える。全列が二直線列の場合は各位相一中心で
系\ref{cor:four-two-line-impossible-with-pure}が適用される。残る全行を示す。
位相番号を揃えるためだけに巡回しており、添字移動で対象を同一視していない。
\[
\begin{array}{c|l}
\text{順中心の位相（列順）}&\text{列の次数列}\\\hline
1,2,3&11,11,131\\
1,2,3&121,12321,131\\
0,1,2,3&121,121,121,121\\
0,1,2,3&1221,11,1221,11
\end{array}
\]
第一行は定理\ref{thm:cubic-corridor-jets}に反する。
第二行は位相1の順・逆中心で始終する $121$ を持つ。位相3の $131$ の中央三次曲線上の
順・逆欠損は各5で、位相1の各向きの唯一の中心に集中する。第三行も各 $121$ の中央二次曲線上の
両欠損が各2である。いずれも埋込みなので両端で過剰となる $121$ があって
系\ref{cor:conic-corridor-consecutive}に反する。第四行は位相0の $1221$ の始点に、
位相3の二直線列の初回直線から順欠損2が集中し、同じ系に反する。
\end{proof}

\begin{corollary}\label{cor:all-positive-dimensional-window-decision}
標準正則窓の四辺は全て純テンソル関係を含まない。また少なくとも一射影が正次元の
標準有限Frobenius候補で、全次数の標準Koszul左右AS正則性と、正則な窓間の添字保持同型は有限に判定できる。
\end{corollary}
\begin{proof}
純関係の不在は本定理と定理\ref{thm:single-pure-branch-exclusion}冒頭の曲面・一点グラフ・有限辺の議論による。
従って純関係をSegre交叉で検査して非空なら排除でき、空なら系\ref{cor:positive-dimensional-pure-free-decision}を使う。
\end{proof}
